% !TEX root = ../../main.tex

\section{Corpus de validación}
\label{sec:validacion-corpus}

La evaluación se organiza en dos corpus sintéticos positivos y un control negativo auxiliar. El primer corpus usa la mónada \texttt{Maybe} para estudiar el comportamiento estructural de la extensión en un contexto determinista. El segundo emplea la mónada probabilística \texttt{Dist.T Rational} para evaluar su aplicación en el dominio principal de la memoria: los programas probabilísticos escritos en \textsf{Haskell}. En conjunto reúnen 99 programas concretos: 47 basados en \texttt{Maybe} y 52 en distribuciones probabilísticas finitas. Su objetivo no es representar exhaustivamente todos los programas posibles, sino cubrir de manera sistemática las formas de dependencia, patrones de binders, clases de statements, comportamientos de costo y fenómenos semánticos que inciden directamente en la solución.

Cada programa se almacena de acuerdo con una jerarquía de familia conceptual, caso específico y variante interna. La variante identifica normalmente la cantidad de statements del bloque \texttt{do} del programa a validar, por ejemplo \texttt{03} o \texttt{04}; cuando existen varias variantes con la misma cantidad de statements se agrega un sufijo descriptivo, como \texttt{08-reorder-improves-ado}.

\subsection{Familias compartidas}

Los corpus comparten 42 variantes estructurales, adaptadas a las operaciones de cada mónada sin modificar la forma de dependencia que se pretende estudiar. La Tabla~\ref{tab:validacion-familias-compartidas} resume su distribución. Como cada variante aparece una vez con \texttt{Maybe} y una vez con \texttt{Dist.T Rational}, este núcleo compartido aporta 84 de los 99 programas.

\begin{table}[ht]
\centering
\TableSetup
\TableFrame{%
\begin{tabular}{p{4.1cm}cp{8.1cm}}
\rowcolor{tableHeaderBg}
\TableHead{Familia} & \TableHead{Variantes} & \TableHead{Cobertura principal} \\
\texttt{dependency-shapes} & 14 & Distintas formas del grafo de precedencia RAW: cadenas, grafo sin aristas, diamantes, fan-in y fan-out. \\
\texttt{binders} & 11 & Tuplas, listas, patrones anidados, wildcards, as-patterns y patrones lazy. \\
\texttt{let-statements} & 4 & Dependencias entre statements \texttt{let} y binders monádicos. \\
\texttt{body-stmts} & 2 & Statements sin binder, independientes y dependientes. \\
\texttt{shadowing} & 4 & Reutilización de nombres fuente con identidades renombradas distintas. \\
\texttt{cost-selection} & 4 & Empates de costo, mínimos únicos y órdenes originales mejorables. \\
\texttt{controls} & 3 & Activación e inactividad de la ruta conmutativa de la extensión. \\
\end{tabular}%
}
\caption{Familias y variantes compartidas por los dos corpus.}
\label{tab:validacion-familias-compartidas}
\end{table}

Estas familias aíslan el comportamiento del Renamer, la recolección de binders, el cálculo de variables libres locales, la construcción del grafo de precedencia RAW y la política de selección. Las variantes de \texttt{cost-selection} se retoman con código probabilístico en la Sección~\ref{sec:resultados-costos}; las formas fan-in y fan-out permanecen dentro de la cobertura del corpus, aunque no se desarrollan como ejemplos visuales independientes.

\subsection{Familias específicas}

La familia \texttt{maybe-failure} aporta cinco variantes específicas al corpus de \texttt{Maybe}. Por su parte, \texttt{probability-distributions} aporta ocho variantes al corpus probabilístico y \texttt{probability-conditioning} agrega otras dos. La Tabla~\ref{tab:validacion-familias-especificas} resume esta distribución adicional al núcleo compartido.

\begin{table}[ht]
\centering
\TableSetup
\TableFrame{%
\begin{tabular}{p{5.2cm}cp{7.2cm}}
\rowcolor{tableHeaderBg}
\TableHead{Familia} & \TableHead{Variantes} & \TableHead{Cobertura principal} \\
\texttt{maybe-failure} & 5 & \texttt{Nothing}, fallos de patrón y acciones que modelan guardas. \\
{\footnotesize\texttt{probability-distributions}} & 8 & Pesos, soportes dependientes, resultados repetidos y constructores de distribuciones. \\
{\footnotesize\texttt{probability-conditioning}} & 2 & Distribuciones imposibles y filtrado explícito de resultados. \\
\end{tabular}%
}
\caption{Familias específicas de cada corpus.}
\label{tab:validacion-familias-especificas}
\end{table}

El corpus de \texttt{Maybe} entrega programas pequeños y deterministas para observar fallos y propiedades estructurales. El corpus de \texttt{Dist.T Rational}, por su parte, conecta la extensión con el dominio principal de la memoria. Todos sus programas imprimen una salida canónica mediante \texttt{Dist.norm}, que agrupa resultados iguales y suma sus probabilidades. La comparación continúa siendo byte a byte, pero se realiza sobre distribuciones normalizadas para no confundir una diferencia de representación interna con una diferencia semántica.

\subsection{Control negativo con IO}

El caso con \texttt{IO} no se contabiliza como un tercer corpus. Se utiliza para mostrar por qué la preservación de dependencias de un reordenamiento válido no basta para una mónada arbitraria y por qué su equivalencia semántica exige una hipótesis explícita de conmutatividad. En \texttt{IO}, dos acciones que no comparten variables pueden tener efectos observables cuyo orden importa. Por ejemplo, dos impresiones independientes pueden producir salidas distintas si se intercambian.

El programa de control declara deliberadamente una instancia \texttt{CommutativeMonad IO} que no satisface la propiedad asumida y utiliza \texttt{CD.do} para activar la extensión. Por tanto, permite comprobar que el marcador constituye un contrato no verificado: si el programador declara como conmutativos efectos que no lo son, un reordenamiento estructuralmente válido puede cambiar la salida.

El control negativo también muestra por qué el compilador no debe inferir conmutatividad a partir de la sola ausencia de dependencias RAW. La extensión explora órdenes alternativos únicamente en bloques marcados explícitamente; para bloques \texttt{do} ordinarios, \texttt{ApplicativeDo} conserva su comportamiento y no se habilita la ruta conmutativa.
